\scanpage{7}
\mspage{3}
2) \uline{Foncteur ess. surjectif} \deleted{\unsure{est}} \inserted{\uline{si} $S$ \uline{géom. unibranche en les points de} $S-U$ \unsure{(ou|(} \inserted{[plus gén, si, posant $S''=S'\times_S S'$, $U'=U\times_S S'$, $U''=U\times_S S''$, $U'\deleted{\unsure{\text{est dense}}}$ \inserted{$(U'')$} rencontre tte composante connexe de $S'$ $(S'')$]} et si \unsure{le normalisé|la normalisée|la normalisation} $S'$ de $S_{\mathrm{red}}$ est \unsure{fini|finie} sur $S$.}

\deleted{a) Supposons}

Partons de $A_U,M,\varphi$, on veut construire $A$. On peut supposer $S$ affine, $S=\operatorname{Spec}\Lambda$. \deleted{Pour \unsure{descendre|avoir|descente} $A$, il suffit de traiter le cas}

a) Réduction au cas $S$ réduit.

Soit. $S_0=S_{\mathrm{réd}}$, $A_0$ schéma abélien sur $S_0$ qui y résout le pb. Par Serre-Tate, on trouve \deleted{des} \deleted{\unsure{un}} schémas abéliens sur les $T_n$, donc un schéma \inserted{formel} abélien \inserted{$A'$} sur \deleted{\unsure{\ell}} $S'=\operatorname{Spec}\widehat\Lambda$. Son image inverse sur $S'_0$ ``est'' un vrai schéma abélien, \deleted{sur $\unsure{S_0}$} savoir l'image inverse de $A_0$ sur $S'_0$.
\diagram{\begin{tikzcd}[ampersand replacement=\&,row sep=small,column sep=small]
 S_0\arrow[d] \& S'_0\arrow[l]\arrow[d]\\
 S \& S'\arrow[l]
\end{tikzcd}}
Donc (ref \ldots\ !!!) $A'$ provient d'un vrai schéma abélien $A'$ sur $S'$ [l'imm. $S'_0-S'$ étant \unsure{nilpotente}]. Or $U\amalg S'\longrightarrow S$ est fid. plat, et on trouve une donnée de descente sur $A_U\amalg A'$. Elle devient effective après chgt de base $S_0\unsure{\longleftarrow|-|\longrightarrow}S$, donc elle l'était déjà avant, d'après la \unsure{section|critère|notion|réduction} de Raynaud. On gagne !

\margin{?}
\deleted{b) Réduction au cas $S$ local, \deleted{\unsure{et pt}} $T$ son point fermé. Pour faire cette réduction, on est ramené à prouver ceci : si $t$ est un pt \unsure{fermé\ldots?} de $T$, $V$ un ouvert \unsure{cont|contenant} $t$ et $U$, \deleted{et} $A_V$ un schéma abélien sur $V$ \deleted{qui résout} tel que sa restriction à $S'=\operatorname{Spec}\mathcal O_{S,t}$ \deleted{résolve} soit OK., alors $A_V$ est \deleted{\unsure{encore}} OK dans un voisinage, i.e. l'isom donné $T_p(A_{S'})\simeq M_{S'}$ se prolonge en un isom en un}

% archive sheet number: 6
